% Based on templates/lecture-example.tex.
% Related lecture: SC2005-L12, 2026-09-18.
% Source examples: Part5 (Deadlocks).pdf, 5.26, 5.31--5.33, 5.36--5.37.
% The last P4 question on 5.37 has no worked solution; the derivation is supplied.

\exercisemeta
  {SC2005-L12-EX}
  {2026-09-18}
  {Part 5，pp.~5.26、5.31--5.33、5.36--5.37（PDF pp.~26、31--33、36--37）}
  {讲义数据与完成序列已核对；p.~5.37 最后一问补全推导}
  {原始讲义与解答已核对}

\exercisequestion{SC2005-L12-E01}{安全序列：存在一个即可}

\textbf{来源：}讲义 p.~5.26，Example of Safe Process Sequence。\quad
\textbf{对应：}\relatedtopic{SC2005-L12-T01}{Safe State 与资源模型}。

只有一种资源，$Available=1$。$P_1,P_2,P_3$ 各持有 $1$ 个实例，表中的 Request 分别为 $1,2,3$，在本例中表示完成前需补齐的数量。判断 $\langle P_1,P_2,P_3\rangle$ 与 $\langle P_3,P_2,P_1\rangle$ 是否安全，并判断整个系统是否 safe。

\begin{checkedsolution}
第一种顺序逐步满足所需数量，过程为
\begin{center}
\begin{tabular}{@{}ccccc@{}}
  \toprule
  进程 & 完成前 Work & 补齐数量 & 原已持有 & 完成后 Work \\
  \midrule
  $P_1$ & $1$ & $1$ & $1$ & $1-1+(1+1)=2$ \\
  $P_2$ & $2$ & $2$ & $1$ & $2-2+(1+2)=3$ \\
  $P_3$ & $3$ & $3$ & $1$ & $3-3+(1+3)=4$ \\
  \bottomrule
\end{tabular}
\end{center}
所以 $\langle P_1,P_2,P_3\rangle$ 安全。第二种顺序一开始要求 $3$ 个额外实例，而当前只有 $1$ 个，因此该顺序不安全。但系统仍 safe，因为已找到至少一个安全序列。一个失败的排列不能证明所有排列都失败。
\end{checkedsolution}

\exercisequestion{SC2005-L12-E02}{五进程、三类资源的 Safety Algorithm}

\textbf{来源：}讲义 pp.~5.31--5.33。\quad
\textbf{对应：}\relatedtopic{SC2005-L12-T02}{Safety Algorithm}。

有进程 $P_0,\ldots,P_4$，资源总量 $Total=(10,5,7)$，当前 $Available=(3,3,2)$。所有三元组均按 $A,B,C$ 排列。下表前两列为题设，计算出的 Need 一并列出，供逐行核对。判断当前状态是否 safe。

\begin{center}
\begin{tabular}{@{}cccc@{}}
  \toprule
  进程 & Allocation & Max & $Need=Max-Allocation$ \\
  \midrule
  $P_0$ & $(0,1,0)$ & $(7,5,3)$ & $(7,4,3)$ \\
  $P_1$ & $(2,0,0)$ & $(3,2,2)$ & $(1,2,2)$ \\
  $P_2$ & $(3,0,2)$ & $(9,0,2)$ & $(6,0,0)$ \\
  $P_3$ & $(2,1,1)$ & $(2,2,2)$ & $(0,1,1)$ \\
  $P_4$ & $(0,0,2)$ & $(4,3,3)$ & $(4,3,1)$ \\
  \bottomrule
\end{tabular}
\end{center}

\begin{checkedsolution}
例如 $Need_1=(3,2,2)-(2,0,0)=(1,2,2)$。初始化 $Work=(3,3,2)$，所有 Finish 为 false；初始 $P_1$ 与 $P_3$ 均可选。按讲义的顺序选择，每行均先验证 $Need_i\le Work$，再加回该进程原有的 Allocation：
\begin{center}
\begin{tabular}{@{}ccccc@{}}
  \toprule
  选择 & 完成前 Work & Need & 加回 Allocation & 完成后 Work \\
  \midrule
  $P_1$ & $(3,3,2)$ & $(1,2,2)$ & $(2,0,0)$ & $(5,3,2)$ \\
  $P_3$ & $(5,3,2)$ & $(0,1,1)$ & $(2,1,1)$ & $(7,4,3)$ \\
  $P_4$ & $(7,4,3)$ & $(4,3,1)$ & $(0,0,2)$ & $(7,4,5)$ \\
  $P_2$ & $(7,4,5)$ & $(6,0,0)$ & $(3,0,2)$ & $(10,4,7)$ \\
  $P_0$ & $(10,4,7)$ & $(7,4,3)$ & $(0,1,0)$ & $(10,5,7)$ \\
  \bottomrule
\end{tabular}
\end{center}
因此 $\langle P_1,P_3,P_4,P_2,P_0\rangle$ 是安全序列，系统 safe。最后 Work 等于 Total，表示模拟中所有资源均已归还；实际的 Available 仍为 $(3,3,2)$，不能被最终 Work 覆盖。
\end{checkedsolution}

\exercisequestion{SC2005-L12-E03}{批准请求：\texorpdfstring{$P_1$ 申请 $(1,0,2)$}{P1 申请 (1,0,2)}}

\textbf{来源：}讲义 p.~5.36。\quad
\textbf{对应：}\relatedtopic{SC2005-L12-T03}{Resource-Request Algorithm}。

从\relatedexercise{SC2005-L12-E02}{五进程的初始快照}出发，$P_1$ 当前持有 $(2,0,0)$，剩余最大需求为 $(1,2,2)$；它申请 $(1,0,2)$。是否可以批准？

\begin{checkedsolution}
上限与库存检查分别为
\[
  (1,0,2)\le(1,2,2),\qquad (1,0,2)\le(3,3,2),
\]
均成立。试分配后
\[
  Available'=(2,3,0),\quad Allocation_1'=(3,0,2),\quad Need_1'=(0,2,0).
\]
其余进程的 Allocation、Need 以及全部 Max 不变。用新的 $Work=(2,3,0)$ 检查讲义给出的序列：
\begin{center}
\begin{tabular}{@{}ccccc@{}}
  \toprule
  选择 & 完成前 Work & Need & 加回 Allocation & 完成后 Work \\
  \midrule
  $P_1$ & $(2,3,0)$ & $(0,2,0)$ & $(3,0,2)$ & $(5,3,2)$ \\
  $P_3$ & $(5,3,2)$ & $(0,1,1)$ & $(2,1,1)$ & $(7,4,3)$ \\
  $P_4$ & $(7,4,3)$ & $(4,3,1)$ & $(0,0,2)$ & $(7,4,5)$ \\
  $P_0$ & $(7,4,5)$ & $(7,4,3)$ & $(0,1,0)$ & $(7,5,5)$ \\
  $P_2$ & $(7,5,5)$ & $(6,0,0)$ & $(3,0,2)$ & $(10,5,7)$ \\
  \bottomrule
\end{tabular}
\end{center}
每一步都满足 $Need_i\le Work$，因此 $\langle P_1,P_3,P_4,P_0,P_2\rangle$ 安全，请求获批。虽然试分配后 C 的空闲数为 $0$，$P_1$ 却不再需要额外 C，可以先完成并释放它持有的 C。

\textbf{后续请求的真实起点：}批准仅落实那一笔 $(1,0,2)$ 的分配，故实际 Available 是 $(2,3,0)$，不是安全模拟最后的 $(10,5,7)$；也没有进程因这次检查而实际结束。
\end{checkedsolution}

\par\noindent\begin{minipage}{\textwidth}
\exercisequestion{SC2005-L12-E04}{后续请求：报错、试分配失败与回滚}

\textbf{来源：}讲义 p.~5.37 的三个后续请求；最后一个 $P_4$ 问题的解答为补充推导。\quad
\textbf{对应：}\relatedtopic{SC2005-L12-T03}{Resource-Request Algorithm}。

共同起点为上一例中 $P_1$ 的请求刚获批、尚无进程实际释放资源的状态，$Available=(2,3,0)$：
\end{minipage}\par
\begin{center}
\begin{tabular}{@{}ccc@{}}
  \toprule
  进程 & Allocation & Need \\
  \midrule
  $P_0$ & $(0,1,0)$ & $(7,4,3)$ \\
  $P_1$ & $(3,0,2)$ & $(0,2,0)$ \\
  $P_2$ & $(3,0,2)$ & $(6,0,0)$ \\
  $P_3$ & $(2,1,1)$ & $(0,1,1)$ \\
  $P_4$ & $(0,0,2)$ & $(4,3,1)$ \\
  \bottomrule
\end{tabular}
\end{center}
分别判断 $P_1$ 再申请 $(1,0,2)$、$P_0$ 申请 $(0,2,0)$、$P_4$ 申请 $(2,3,0)$。失败请求不改变上述共同起点。

\begin{checkedsolution}
\textbf{(a) $P_1$ 再申请 $(1,0,2)$。}该向量不满足 $Request_1\le Need_1=(0,2,0)$，A、C 均超出剩余上限，应报错，不进入试分配。

\textbf{(b) $P_0$ 申请 $(0,2,0)$。}上限检查 $(0,2,0)\le(7,4,3)$ 与库存检查 $(0,2,0)\le(2,3,0)$ 均通过。试分配后
\[
  Available'=(2,1,0),\quad Allocation_0'=(0,3,0),\quad Need_0'=(7,2,3).
\]
以 $Work=(2,1,0)$ 尝试寻找第一个可完成的进程：
\begin{center}
\begin{tabular}{@{}ccl@{}}
  \toprule
  进程 & 试分配后的 Need & 无法满足的资源类型 \\
  \midrule
  $P_0$ & $(7,2,3)$ & A、B、C \\
  $P_1$ & $(0,2,0)$ & B：需 $2$，仅有 $1$ \\
  $P_2$ & $(6,0,0)$ & A：需 $6$，仅有 $2$ \\
  $P_3$ & $(0,1,1)$ & C：需 $1$，仅有 $0$ \\
  $P_4$ & $(4,3,1)$ & A、B、C \\
  \bottomrule
\end{tabular}
\end{center}
没有任何候选，故试算状态 unsafe。撤销本次试分配，使 Available 恢复 $(2,3,0)$、$Allocation_0$ 恢复 $(0,1,0)$、$Need_0$ 恢复 $(7,4,3)$，让 $P_0$ 等待。讲义红字中的严格小于应按正式算法取 $Need_i\le Work$；相等也能满足。

\textbf{(c) $P_4$ 申请 $(2,3,0)$：补充推导。}从恢复后的共同起点出发，$(2,3,0)\le(4,3,1)$ 且 $(2,3,0)\le(2,3,0)$。试分配后
\[
  Available'=(0,0,0),\quad Allocation_4'=(2,3,2),\quad Need_4'=(2,0,1).
\]
此时各进程的 Need 都不是零向量，因而没有任何 $Need_i\le Work=\mathbf0$ 的候选。状态 unsafe，应恢复共同起点并让 $P_4$ 等待。结论依赖“没有零 Need 进程”，不能仅凭 Available 为零就一概判为 unsafe。

三问中，(a) 是超过声明上限的错误；(b)、(c) 是合法且库存足够、但安全检查不通过的请求。后两者要回滚并等待，不能把试算 unsafe 直接写成已经检测到实际死锁。
\end{checkedsolution}
